\documentclass[tikz,border=4pt]{standalone}
\usepackage{ctex}
\usepackage{amsmath}
\usetikzlibrary{calc, arrows.meta, decorations.markings}

\begin{document}
\begin{tikzpicture}[scale=1.1, thick]
  % 说明一元一次方程 ax+b=0 的结构

  % 天平底座
  \draw[black, line width=1.5pt] (-3.5,0) -- (3.5,0);
  \draw[black, line width=1.5pt] (0,0) -- (0,-0.6);
  \draw[black, line width=1.5pt] (-0.5,-0.6) -- (0.5,-0.6);

  % 天平横梁
  \draw[black, line width=1.5pt] (-3,0.8) -- (3,0.8);
  \draw[black] (0,0) -- (0,0.8);

  % 左盘（含 ax+b）
  \draw[gray, dashed] (-3,0.8) -- (-3,0.3);
  \draw[black] (-3.8,0.3) -- (-2.2,0.3);
  \node[above, blue] at (-3,0.6) {$ax + b$};

  % 右盘（含 0）
  \draw[gray, dashed] (3,0.8) -- (3,0.3);
  \draw[black] (2.2,0.3) -- (3.8,0.3);
  \node[above, blue] at (3,0.6) {$0$};

  % 等号强调
  \node[red, font=\Large\bfseries] at (0,1.4) {$=$};

  % 标注
  \node[below] at (0,-1.0) {\small 等式两边保持平衡——天平的形象};
  \node[below] at (0,-1.45) {\small 性质1：两边同加减同数，等式不变};
  \node[below] at (0,-1.85) {\small 性质2：两边同乘除同非零数，等式不变};
\end{tikzpicture}
\end{document}
